\documentclass[11pt]{article}
\usepackage[T1]{fontenc}
\usepackage[utf8]{inputenc}
\usepackage{lmodern}
\usepackage[margin=1in]{geometry}
\usepackage{amsmath,amssymb,amsthm,mathtools}
\usepackage{booktabs,array,longtable}
\usepackage{microtype}
\usepackage{enumitem}
\usepackage{needspace}
\usepackage{xcolor}
\usepackage[colorlinks=true,linkcolor=blue!35!black,citecolor=blue!35!black,urlcolor=blue!35!black]{hyperref}
\usepackage{fancyvrb}
\usepackage[capitalise,noabbrev]{cleveref}
\newcommand{\AuthorName}{Ziheng Wang}
\newcommand{\AuthorAffiliation}{University of Chinese Academy of Sciences,\\Hangzhou Institute for Advanced Study, Hangzhou, China}

\newtheorem{theorem}{Theorem}[section]
\newtheorem{proposition}[theorem]{Proposition}
\newtheorem{lemma}[theorem]{Lemma}
\newtheorem{corollary}[theorem]{Corollary}
\theoremstyle{definition}
\newtheorem{definition}[theorem]{Definition}
\theoremstyle{remark}
\newtheorem{remark}[theorem]{Remark}
\AddToHook{env/theorem/begin}{\crefalias{theorem}{theorem}}
\AddToHook{env/proposition/begin}{\crefalias{theorem}{proposition}}
\AddToHook{env/lemma/begin}{\crefalias{theorem}{lemma}}
\AddToHook{env/corollary/begin}{\crefalias{theorem}{corollary}}
\numberwithin{equation}{section}
\newcommand{\Z}{\mathbb Z}
\newcommand{\Q}{\mathbb Q}
\newcommand{\C}{\mathbb C}
\newcommand{\SL}{\mathrm{SL}}
\newcommand{\F}{F(a,b)}
\newcommand{\low}[1]{\begin{pmatrix}1&0\\#1&1\end{pmatrix}}
\newcommand{\up}[1]{\begin{pmatrix}1&#1\\0&1\end{pmatrix}}
\newcommand{\ev}{\operatorname{ev}}
\newcommand{\syll}{\lVert\,\cdot\,\rVert_{\mathrm{syl}}}
\setlist{itemsep=3pt,topsep=5pt}
\setlength{\emergencystretch}{2em}
\ifdefined\pdfinfoomitdate\pdfinfoomitdate=1\pdftrailerid{}\pdfsuppressptexinfo=15\fi
\hypersetup{pdftitle={Nonfreeness for rational two-parabolic groups with numerator 24},pdfauthor={\AuthorName}}
\title{Nonfreeness for rational two-parabolic groups\protect\\with numerator 24}
\author{\AuthorName\ifx\AuthorAffiliation\empty\else\\\small\AuthorAffiliation\fi}
\date{October 10, 2026}
\begin{document}
\maketitle
\begin{abstract}
Let $A=\left(\begin{smallmatrix}1&0\\1&1\end{smallmatrix}\right)$ and
$B_q=\left(\begin{smallmatrix}1&q\\0&1\end{smallmatrix}\right)$.
We prove that $\langle A,B_{24/r}\rangle\leq\SL_2(\C)$ is not a free group
for every integer $r$ with $|r|>6$. Each parameter has a nontrivial
relation of at most $64$ syllables. Moreover, $24/r$ is a two-step
relation number whenever $|r|>17$. The proof combines a finite
certificate of $48$ good residue-class rules with a five-syllable
identity for every nonzero integer $x$ and $R=1+1260x$:
\[
 B_{24/R}^{\,3Rx} A^{18+630x} B_{24/R}^{-3}
 A^{18x}B_{24/R}^{\,R(3+105x)}
 =\begin{pmatrix}1&0\\18/R&1\end{pmatrix}.
\]
This identity covers the four residue classes left by the finite
certificate. All finite data and their exact verification are included,
and the accompanying Lean project proves the universal nonfreeness
and bounded-relation statements.
\end{abstract}

\section{Introduction}
For $q\in\C$, put
\begin{equation}\label{eq:generators}
 A=\low{1},\qquad B_q=\up{q},\qquad
 G_q=\langle A,B_q\rangle\leq\SL_2(\C).
\end{equation}
The rational two-parabolic problem asks whether $G_q$ is nonfree for
every nonzero rational $q$ with $|q|<4$. Kim and Koberda
\cite[Theorem~1.4]{KK2020} proved this for parameters $q=s/r$ with
$0<|s|\leq27$ and $|s|\ne24$. For numerator $24$, they proved that
the denominators giving two-step relation numbers have natural density
one. We prove the all-denominator conclusion for this numerator.

A \emph{syllable} is a nonzero integer power of one formal generator.
The syllable length of a freely reduced word counts its alternating
power blocks; it does not count the letters obtained by expanding
those powers. Let $\F$ denote the free group on $a,b$, and let
$\ev_q:\F\to\SL_2(\C)$ send $a$ to $A$ and $b$ to $B_q$.

\begin{theorem}\label{thm:main}
For every integer $r$ satisfying $|r|>6$, the group $G_{24/r}$ is not
an abstract free group. There is a nonempty freely reduced word
$v_r\in\F$ with at most $64$ syllables such that
\[
 v_r\ne1,\qquad \ev_{24/r}(v_r)=I.
\]
No coprimality assumption on $r$ and $24$ is required.
\end{theorem}

Following \cite[Definition in Section~1]{KK2020}, a nonzero parameter
$q$ is a \emph{two-step relation number} if an alternating word
\begin{equation}\label{eq:twostep-def}
 b^{e_1}a^{e_2}\cdots b^{e_{2k+1}},\qquad
 0\leq k\leq2,\qquad e_j\in\Z\setminus\{0\},
\end{equation}
has a lower-triangular value under $\ev_q$.

\begin{corollary}\label{cor:twostep}
The parameter $24/r$ is a two-step relation number for every integer
$r$ with $|r|>17$. In particular, the set of such denominators is
cofinite in $\Z\setminus\{0\}$.
\end{corollary}

\begin{corollary}\label{cor:smallnumerators}
If $s,r\in\Z$, $r\ne0$, $0<|s|\leq27$, and $|s/r|<4$, then
$G_{s/r}$ is not a free group.
\end{corollary}
\begin{proof}
Use \cref{thm:main} when $|s|=24$, and
\cite[Theorem~1.4(1)]{KK2020} otherwise. The identity
$B_{-q}=B_q^{-1}$ handles the signs.
\end{proof}
The last corollary uses the existing results for other numerators.
The argument below is self-contained for numerator $24$; it makes no
claim for all rational parameters in $(-4,4)$.

\subsection*{The mechanism of the proof}
The good-class criterion of Kim and Koberda reduces many parameters
to a five-syllable lower-unipotent matrix. We use $48$ explicit rules,
including rules for denominators that are not coprime to $24$.
Together they cover $40316$ of the $40320$ residue classes modulo
$40320$. The missing representatives are
\begin{equation}\label{eq:omitted}
 1261,\quad6299,\quad34021,\quad39059.
\end{equation}
Each is congruent to $1$ or $-1$ modulo $1260$. The same
five-syllable mechanism admits a particularly simple specialization
on these two progressions: set $R=1+1260x$ and use
\[
 m=18x,\quad y=3x,\quad u=1+36x,\quad t=3,\quad e=-3-105x.
\]
Then $e=y-ut$ and $Ru+24me=1$. These two equations give the
identity in the abstract, with no recursive descent or valuation
calculation. The finitely many exceptions to the good-class rules
are handled by $23$ exact words.

The logical roles of the ingredients are separate. The residue-class
check is a complete finite computation, the parametric identity
handles every integer in the remaining infinite families, and free
reduction proves that the resulting relations are nontrivial.
We also give the group-theoretic argument that a nontrivial relation
for these two noncommuting generators implies abstract nonfreeness.
Section~\ref{sec:formal} identifies the corresponding Lean statements
and the verification boundary.

\section{A five-syllable matrix identity}\label{sec:identity}
Write $L(z)=\left(\begin{smallmatrix}1&0\\z&1\end{smallmatrix}\right)$
and $U(z)=\left(\begin{smallmatrix}1&z\\0&1\end{smallmatrix}\right)$.
For every integer $n$,
\begin{equation}\label{eq:powers}
 A^n=L(n),\qquad B_q^n=U(nq).
\end{equation}
Our commutator convention is $[X,Y]=XYX^{-1}Y^{-1}$.

\Needspace{11\baselineskip}
\begin{lemma}\label{lem:five}
Let $s,r,m,y,u,t,e,k\in\Z$, with $r\ne0$, satisfy
\begin{equation}\label{eq:five-assumptions}
 m=ky,\qquad e=y-ut,\qquad ru+sme=1.
\end{equation}
For $B=B_{s/r}$, one has
\begin{equation}\label{eq:five}
 W:=B^{ry}A^{-ke}B^{-t}A^mB^{-re}=L(kt/r).
\end{equation}
If $m,y,t,e$ are all nonzero, the formal word $w$ displayed in
\eqref{eq:five} has five nonzero syllables, and $[w,a]$ is a
nontrivial relation of exactly $12$ syllables.
\end{lemma}
\begin{proof}
Equations \eqref{eq:powers} and \eqref{eq:five-assumptions} give
\begin{align*}
 B^{ry}A^{-ke}
 &=\begin{pmatrix}1-skye&sy\\-ke&1\end{pmatrix}
  =\begin{pmatrix}ru&sy\\-ke&1\end{pmatrix},\\
 B^{ry}A^{-ke}B^{-t}
 &=\begin{pmatrix}ru&se\\-ke&1+sket/r\end{pmatrix}.
\end{align*}
After multiplication by $A^m$, the upper-left entry is $ru+sme=1$.
If $K$ is the lower-left entry, then
\begin{align*}
 rK&=rm-rke+skemt\\
   &=k\bigl(r(y-e)+smet\bigr)\\
   &=kt(ru+sme)=kt.
\end{align*}
Thus
\[
 B^{ry}A^{-ke}B^{-t}A^m
 =\begin{pmatrix}1&se\\kt/r&1+sket/r\end{pmatrix}.
\]
Right multiplication by $B^{-re}=U(-se)$ proves \eqref{eq:five}.
The matrices $L(kt/r)$ and $A$ commute.

If $m,y,t,e$ are nonzero, then $k\ne0$, and all five exponents in
$w=b^{ry}a^{-ke}b^{-t}a^mb^{-re}$ are nonzero. In
$[w,a]=waw^{-1}a^{-1}$, the generator types alternate as
$b,a,b,a,b,a,b,a,b,a,b,a$. Consequently the word is freely
reduced and nonempty, while its matrix value is $I$.
\end{proof}

The next lemma is the good-class criterion
\cite[Definition~3.3 and Lemma~3.5]{KK2020}, stated with the
explicit relation length needed here. We include the proof, in
particular its normalization when the denominator is not coprime
to the numerator.

\begin{lemma}[Good residue classes]\label{lem:good}
Let $s,w,m,y$ be nonzero integers, and put
\[
 D=\gcd(w,sm)>0,\qquad d=\gcd(w,s)>0.
\]
Suppose
\begin{equation}\label{eq:goodconditions}
 yD\mid md,\qquad w\mid smy-D\ \text{or}\ w\mid smy+D.
\end{equation}
Every integer
\[
 r\equiv w\pmod{sm},\qquad r\notin\{0,w,D,-D\},
\]
admits a five-syllable lower-unipotent word for $s/r$, and therefore
a nontrivial $12$-syllable relation.
\end{lemma}
\begin{proof}
First consider the normalized situation
\[
 y\mid m,\qquad w\mid smy-1,\qquad
 r=w+smt\notin\{0,w,1,-1\}.
\]
Set $u=(1-smy)/w$, choose $k\in\Z$ with $m=ky$, and put
$e=y-ut$. Then
\[
 ru+sme=(w+smt)u+sm(y-ut)=wu+smy=1.
\]
Here $t\ne0$ because $r\ne w$. Also $e\ne0$, since $e=0$
would give $ru=1$ and hence $r\in\{1,-1\}$. Apply
\cref{lem:five}. If instead $w\mid smy+1$, replace $y$ by $-y$.
This changes neither divisibility by $y$ nor its nonvanishing.

For the general case, put
\[
 c=D/d,\quad w'=w/D,\quad r'=r/D,\quad
 s'=s/d,\quad m'=md/D.
\]
All these quantities are integers. Indeed,
$D/d=\gcd(w/d,m)$ because $\gcd(w/d,s/d)=1$, so $c\mid m$.
Moreover, $D\mid r$ and
\[
 r'=w'+s'm't,\qquad y\mid m',\qquad
 w'\mid s'm'y-1\ \text{or}\ w'\mid s'm'y+1,
\]
where $r=w+smt$ and $t\ne0$. The exclusions on $r$ give
$r'\notin\{0,w',1,-1\}$. The normalized argument produces a
five-syllable word for $s'/r'$, even when $s'=1$.
Finally,
\[
 \frac{s'}{r'}=c\frac sr,\qquad B_{s'/r'}=(B_{s/r})^c.
\]
Multiply every $b$ exponent by the nonzero integer $c$.
The matrix identity, nonzero exponents, and alternating syllable
pattern are preserved.
\end{proof}

\section{The two remaining progressions}\label{sec:direct}
\begin{proposition}\label{prop:direct}
For every nonzero integer $x$, let $R=1+1260x$ and $B=B_{24/R}$.
Then
\begin{equation}\label{eq:direct}
 B^{3Rx}A^{18+630x}B^{-3}A^{18x}B^{R(3+105x)}=L(18/R).
\end{equation}
The commutator of the displayed formal word with $a$ is a
nontrivial $12$-syllable relation.
\end{proposition}
\begin{proof}
Take
\begin{equation}\label{eq:directparameters}
 m=18x,\quad y=3x,\quad k=6,\quad u=1+36x,
 \quad t=3,\quad e=-3-105x.
\end{equation}
Then $m=ky$ and $e=y-ut$. A direct expansion gives
\begin{align*}
 Ru+24me
 &=(1+1260x)(1+36x)+432x(-3-105x)\\
 &=1+1296x+45360x^2-1296x-45360x^2=1.
\end{align*}
Since $x$ is a nonzero integer, $|R|\ge1259$, and $m,y,t$ are
nonzero. The equality $e=0$ would imply $35x=-1$, which is
impossible in $\Z$. Thus \cref{lem:five} applies. Substituting
\eqref{eq:directparameters} into \eqref{eq:five} gives precisely
\eqref{eq:direct}, with $kt=18$.
\end{proof}

\begin{corollary}\label{cor:progressions}
If $r\equiv\pm1\pmod{1260}$ and $|r|>1$, then $24/r$ has a
five-syllable lower-unipotent word and a nontrivial $12$-syllable
relation. In particular, this holds for every integer whose residue
modulo $40320$ is in \eqref{eq:omitted}.
\end{corollary}
\begin{proof}
Choose $\varepsilon\in\{1,-1\}$ such that $R=\varepsilon r$
is congruent to $1$ modulo $1260$. Write $R=1+1260x$.
Then $x\ne0$, since $|r|>1$. Apply \cref{prop:direct}.
Because $B_{24/R}=B_{24/r}^{\varepsilon}$, multiply the three
$b$ exponents in \eqref{eq:direct} by $\varepsilon$. This does
not change nonvanishing or syllable length. For the final assertion,
$1260\mid40320$, and the four representatives in \eqref{eq:omitted}
reduce to $1,-1,1,-1$ modulo $1260$.
\end{proof}

For example, $R=1261$ corresponds to $x=1$, and
\[
 B_{24/1261}^{3783}A^{648}B_{24/1261}^{-3}
 A^{18}B_{24/1261}^{136188}=L(18/1261).
\]
For $r=6299$, choose $R=-6299$ and $x=-5$. In the original
parameter $24/6299$, this becomes
\[
 B_{24/6299}^{-94485}A^{-3132}B_{24/6299}^{3}
 A^{-90}B_{24/6299}^{-3288078}=L(-18/6299).
\]
The examples illustrate the sign convention; the proof of
\cref{prop:direct} applies to all nonzero integers $x$.

\section{The finite certificate}\label{sec:finite}
Let $\mathcal T$ be the $48$ quadruples $(w,m,y,D)$ in
\cref{app:rules}, and let $M=40320$. Each row denotes both residue
classes $r\equiv w$ and $r\equiv-w\pmod{24m}$. Define
\begin{equation}\label{eq:cover}
 \mathcal C=\bigcup_{(w,m,y,D)\in\mathcal T}
 \{a\in\{0,\ldots,M-1\}:a\equiv\pm w\pmod{24m}\}.
\end{equation}
Let
\begin{equation}\label{eq:exceptions}
\begin{split}
 E=\{&7,8,9,10,11,12,13,14,17,24,31,35,37,59,61,83,85,\\
     &107,109,613,877,1139,2267\}.
\end{split}
\end{equation}

\begin{proposition}[Finite covering certificate]\label{prop:finite}
The table $\mathcal T$ has the following properties.
\begin{enumerate}[label=\textup{(\roman*)}]
\item Every row satisfies $24m\mid M$, $D=\gcd(w,24m)$,
$yD\mid m\gcd(w,24)$, and $w\mid24my-D$ or $w\mid24my+D$.
\item The complement of $\mathcal C$ is precisely
$\{1261,6299,34021,39059\}$.
\item The set of numbers greater than $6$ occurring as a $w$ or
$D$ in the table is exactly $E$.
\end{enumerate}
\end{proposition}
\begin{proof}
These are finite integer assertions about the complete table in
\cref{app:rules}. For clarity, the entire covering check is the
following enumeration: start with $\mathcal C=\varnothing$; for each
row and each $0\leq a<M$, insert $a$ if $24m$ divides $a-w$ or
$a+w$; finally take the complement in $\{0,\ldots,M-1\}$.
The result is the four-element set in (ii). The rowwise divisibilities
and the set in (iii) are checked by ordinary integer arithmetic.
The accompanying file \texttt{verify\_certificate.py} performs
these exact checks from \texttt{data/compact\_certificate.json};
\texttt{FiniteCover.lean} proves the same finite assertions using
kernel-checked decision procedures. The tables and the stated tests
fully specify the certificate, independently of how its rows were found.
\end{proof}

\begin{lemma}[Exceptional words]\label{lem:exceptional}
For each $r\in E$, the exponent list in \cref{app:words} defines
an alternating word $w_r=b^{e_1}a^{e_2}\cdots b^{e_L}$ with
nonzero exponents and $L\leq15$. Its value is
\begin{equation}\label{eq:triangularvalue}
 W_r=\begin{pmatrix}\alpha&0\\\beta&\alpha^{-1}\end{pmatrix},
\end{equation}
where the exact rational numbers $\alpha,\beta$ are given in
\cref{app:matrices}. Consequently $[w_raw_r^{-1},a]$ is a
nontrivial relation with $4L+4\leq64$ syllables.
\end{lemma}
\begin{proof}
Multiply the matrices in \eqref{eq:powers} using each exponent list.
For a direct scalar recurrence, if the accumulated matrix is
$\left(\begin{smallmatrix}p&q\\v&z\end{smallmatrix}\right)$,
right multiplication by $B_{24/r}^n$ or $A^n$ respectively gives
\[
 (p,q,v,z)\longmapsto(p,q+24np/r,v,z+24nv/r),
\]
\[
 (p,q,v,z)\longmapsto(p+nq,q,v+nz,z).
\]
Starting from $(1,0,0,1)$ gives exactly the entries in
\cref{app:matrices}. This also specifies a short exact-rational
verification of every row.

For any matrix of the form \eqref{eq:triangularvalue},
\[
 W_rAW_r^{-1}=L(\alpha^{-2}),
\]
which commutes with $A$. The formal relation expands as
\[
 w_raw_r^{-1}a\,w_ra^{-1}w_r^{-1}a^{-1}.
\]
Each copy of $w_r$ or $w_r^{-1}$ begins and ends with a $b$
syllable, and the connecting syllables are nonzero powers of $a$.
There is no free cancellation, so its syllable length is exactly
$4L+4$. The table has $L\leq15$.
\end{proof}

\begin{proposition}\label{prop:allrelations}
For every integer $r$ with $|r|>6$, there is a nontrivial relation
for $A,B_{24/r}$ of at most $64$ syllables.
\end{proposition}
\begin{proof}
Since $B_{-q}=B_q^{-1}$, it suffices to handle $r>6$.
If $r\in E$, use \cref{lem:exceptional}. Otherwise, if its
residue modulo $M$ is one of \eqref{eq:omitted}, use
\cref{cor:progressions}. In the remaining case,
\cref{prop:finite}(ii) supplies a row with
$r\equiv\pm w\pmod{24m}$. Replace $r$ by $-r$ if necessary
to obtain the positive-sign class, apply \cref{lem:good}, and
then invert the $b$ generator back. All excluded values
$0,\pm w,\pm D$ with absolute value greater than $6$ belong to
$\pm E$ by \cref{prop:finite}(iii); our chosen denominator is
therefore not excluded. The good-class and progression constructions
have $12$ syllables, and the exceptional words give at most $64$.
\end{proof}

\begin{proof}[Proof of \cref{cor:twostep}]
Outside $\pm E$, the proof just given produces a five-syllable
lower-unipotent word. For $r\in E$ with $r>17$, every word
listed in \cref{app:words} has three or five syllables. These
are precisely witnesses of the form \eqref{eq:twostep-def}.
Negating the $b$ exponents handles negative denominators.
\end{proof}

\section{From relations to abstract nonfreeness}\label{sec:nonfree}
For completeness we give the group-theoretic bridge, including
the Hopf property used in the formal development. This avoids
identifying failure of freeness of the displayed generating pair
with abstract nonfreeness without justification.

\begin{lemma}\label{lem:hopf}
Every surjective endomorphism of a finitely generated free group
is injective.
\end{lemma}
\begin{proof}
First, a nonidentity element $w$ in any free group is detected by
a finite permutation representation. Choose a word spelling $w$,
and let $S$ be the finite set of its suffix products, including $1$
and $w$. For each free generator $g$, left multiplication by $g$
defines a partial bijection on $S$, on the domain where both the
initial and final elements lie in $S$. Extend it to a permutation
of $S$ by matching the two equally sized complements. The inverse
permutation agrees with multiplication by $g^{-1}$ wherever both
endpoints lie in $S$. By the universal property of the free group,
these permutations define a homomorphism to $\operatorname{Sym}(S)$.
Reading the chosen word from right to left sends $1$ successively
through its suffix products and finally to $w\ne1$. The homomorphism
therefore detects $w$.

Now let $F$ be finitely generated and let $\phi:F\to F$ be
surjective. For any finite group $K$, the set
$\operatorname{Hom}(F,K)$ is finite. Precomposition
$f\mapsto f\circ\phi$ is injective, hence surjective, on this
set. If a nonidentity $w$ belonged to $\ker\phi$, choose a
homomorphism $f:F\to K$ detecting it as in the first paragraph.
Write $f=g\circ\phi$. Then $f(w)=g(1)=1$, a contradiction.
\end{proof}

\begin{lemma}\label{lem:bridge}
Let $H$ be generated by two noncommuting elements $A,B$.
If the evaluation homomorphism $\F\to H$ has nontrivial kernel,
then $H$ is not an abstract free group.
\end{lemma}
\begin{proof}
Suppose that $H$ were free. Its noncommutativity forces a free basis
to contain at least two elements. Sending two distinct basis elements
to $a,b$, and all other basis elements to $1$, gives a surjection
$H\to\F$. Composing this map with the surjective evaluation map
$\F\to H$ gives a surjective endomorphism of $\F$.
By \cref{lem:hopf} it is injective, forcing the evaluation map to
be injective as well. This contradicts its nontrivial kernel.
\end{proof}

\begin{proof}[Proof of \cref{thm:main}]
The relation and its length bound follow from
\cref{prop:allrelations}. For $q=24/r\ne0$,
\[
 AB_q=\begin{pmatrix}1&q\\1&1+q\end{pmatrix},\qquad
 B_qA=\begin{pmatrix}1+q&q\\1&1\end{pmatrix},
\]
so $A$ and $B_q$ do not commute. Apply \cref{lem:bridge}.
All matrices and identities already lie in $\SL_2(\Q)$.
The entrywise inclusion $\SL_2(\Q)\hookrightarrow\SL_2(\C)$
is injective and identifies the generated subgroups, so the same
abstract nonfreeness conclusion holds in $\SL_2(\C)$.
\end{proof}

\section{Formal statements and reproducibility}\label{sec:formal}
The accompanying project is pinned to Lean~4.24.0 and
mathlib~v4.24.0. The source file \texttt{Challenge.lean} specifies
the displayed matrices and target propositions independently of
the implementation; \texttt{Solution.lean} proves those
propositions. The implementation uses the actual free group on a
two-element type, with its standard nontriviality criterion.
A syllable is a pair in $\mathrm{Bool}\times\Z$, with
\texttt{false} denoting $a$ and \texttt{true} denoting $b$.
The reducedness predicate requires alternating generator types and
nonzero exponents. Its connection to a nonidentity free-group element
is proved by expansion into signed letters.

The central bounded-relation theorem has the mathematical content
\begin{align*}
 \forall r\in\Z,\quad |r|>6\ \Longrightarrow\ 
 \exists \ell: &\ell\text{ is nonempty and reduced},\quad |\ell|\leq64,\\
 &\operatorname{word}(\ell)\ne1,\quad
 \ev_{24/r}(\operatorname{word}(\ell))=I.
\end{align*}
The separate nonfreeness theorem concludes
$\neg\operatorname{IsFreeGroup}(G_{24/r})$. The rational-to-complex
map is proved injective, and its image is proved to be exactly the
subgroup generated by the complex versions of \eqref{eq:generators}.
Thus the complex conclusion is an explicit formal theorem.

\begin{center}\small
\begin{tabular}{p{.34\textwidth}p{.58\textwidth}}
\toprule
Mathematical component & Lean source or declaration\\
\midrule
Matrices and reduced words & \texttt{Core.lean}\\[3pt]
General good-class criterion & \texttt{GoodRuleArithmetic.lean}, \texttt{GoodRule.lean}\\[3pt]
Single-parameter identity & \texttt{SimpleDirect.lean}\\[3pt]
Finite cover and exceptional words & \texttt{FiniteCover.lean}, \texttt{FiniteWords.lean}\\[3pt]
Hopf property and nonfreeness & \texttt{FreeGroupHopf.lean}, \texttt{NonFree.lean}\\[3pt]
Uniform syllable bound & \texttt{TwentyFour.twentyfour\_boundedRelation}\\[3pt]
Two-step conclusion for $|r|>17$ & \texttt{TwentyFour.twentyfour\_twoStep}\\[3pt]
Complex matrix-group conclusion & \texttt{TwentyFour.twentyfour\_complex\_not\_free}\\
\bottomrule
\end{tabular}
\end{center}
The single-parameter construction in \cref{prop:direct} is proved in
\texttt{SimpleDirect.lean}; the declaration
\texttt{simple\_direct\_power\_identity} states the five integer powers
in \eqref{eq:direct} directly. Both the universal main proof and the
bounded-relation proof use this construction.

In \texttt{TwoStep.lean}, the theorem \texttt{twentyfour\_twoStep}
separately proves \cref{cor:twostep}. Its predicate requires a
nonzero exponent list of length $1$, $3$, or $5$, beginning with $b$,
whose matrix has zero upper-right entry. This proof checks the
triangular witness directly rather than inferring it from the
$64$-syllable relation bound.

The final source revision passed Lean verification against the
independent specifications, an axiom-dependency audit, and independent
nanoda replay. The replay includes nonfreeness over $\Q$ and $\C$,
nontrivial relation existence, the $64$-syllable bound, the two-step
conclusion, and the direct parametric identity and certificates.
The accompanying verification report records the exact target list
and source and proof-export hashes.


The finite coverage proof checks $40$ blocks of $1008$ residues,
which together exhaust $0,\ldots,40319$. It uses Lean's ordinary
\texttt{decide}, producing proof terms checked by the kernel.
The project does not use \texttt{sorry}, user-added axioms,
\texttt{native\_decide}, or compiler-oracle axioms. The audited
logical dependencies are restricted to Lean's standard
\texttt{propext}, \texttt{Classical.choice}, and \texttt{Quot.sound}.
Independent nanoda replay checks exported declarations and their
transitive dependencies under the same allowed-axiom list.

No acceptance by the full sandboxed Palomar/comparator workflow is
claimed; compatibility and host limitations are documented in the
repository. Version pins, commands, exact theorem lists, and raw
checker logs are supplied there. The standalone Python verifier
uses only integers and \texttt{fractions.Fraction}; it is a convenient
second implementation of the finite arithmetic, rather than a
trusted step in the Lean proof.

\subsection*{AI assistance}
GPT-6-Astra assisted with research exploration, proof construction,
exact computation, Lean formalization, and preparation of this manuscript.

\smallskip
\noindent\copyright\ 2026 Ziheng Wang. This manuscript is licensed under
\href{https://creativecommons.org/licenses/by/4.0/}{Creative Commons Attribution 4.0 International}.

\appendix
\crefalias{section}{appendix}
\clearpage
\section{Complete good-class table}\label{app:rules}
Each row $(w,m,y,D)$ represents the two classes
$r\equiv\pm w\pmod{24m}$. The full arithmetic conditions are
those in \cref{prop:finite}(i). The row numbers are for reference
only; the machine-readable certificate contains the same ordered
quadruples.
\begin{center}\small
\begin{tabular}{rrrrr@{\qquad}rrrrr}
\toprule
$i$&$w$&$m$&$y$&$D$&$i$&$w$&$m$&$y$&$D$\\
\midrule
1&1&1&1&1&25&13&7&1&1\\
2&2&1&1&2&26&35&7&1&7\\
3&3&1&1&3&27&107&7&7&1\\
4&4&1&1&4&28&13&8&4&1\\
5&5&1&1&1&29&59&8&4&1\\
6&6&1&1&6&30&13&10&2&1\\
7&8&1&1&8&31&35&10&1&5\\
8&9&1&1&3&32&37&10&2&1\\
9&12&1&1&12&33&13&15&3&1\\
10&24&1&1&24&34&35&15&3&5\\
11&7&2&1&1&35&83&15&3&1\\
12&10&2&1&2&36&37&16&8&1\\
13&14&2&2&2&37&83&16&8&1\\
14&7&3&3&1&38&59&30&5&1\\
15&31&3&3&1&39&61&30&5&1\\
16&11&4&4&1&40&85&30&2&5\\
17&35&4&4&1&41&37&84&2&1\\
18&7&5&1&1&42&109&84&2&1\\
19&11&5&1&1&43&59&140&20&1\\
20&17&5&1&1&44&1139&140&20&1\\
21&11&6&1&1&45&83&280&28&1\\
22&13&6&1&1&46&2267&280&28&1\\
23&17&6&2&1&47&613&560&40&1\\
24&11&7&7&1&48&877&560&40&1\\
\bottomrule
\end{tabular}
\end{center}


\section{Exact finite verification}\label{app:verification}
Given the table above, the finite part can be checked as follows.
\begin{enumerate}
\item For each row, compute the two gcd values and check the
gcd and divisibility conditions of \cref{prop:finite}(i).
\item Evaluate the predicate in \eqref{eq:cover} for every integer
$a$ from $0$ to $40319$. Its false values must be exactly
$1261,6299,34021,39059$.
\item Collect the values $w,D$ exceeding $6$, remove repetitions,
and compare with \eqref{eq:exceptions}.
\item For every word in \cref{app:words}, apply the rational
recurrences in \cref{lem:exceptional} and compare the resulting
matrix with \cref{app:matrices}.
\end{enumerate}
The supplied verifier also checks the full commutator matrices,
reduced syllable patterns, and the bound $64$. It regenerates all
three tables from the checked data, so the exponent lists and
matrix entries need not be copied manually.

\clearpage
\section{Complete exceptional exponent lists}\label{app:words}
For each listed denominator $r$, put $B=B_{24/r}$. The list
$(e_1,\ldots,e_L)$ denotes
\[
 W_r=B^{e_1}A^{e_2}B^{e_3}\cdots B^{e_L}.
\]
All exponents are nonzero, and every list has odd length. The
relation used uniformly is $[W_rAW_r^{-1},A]=I$, with $4L+4$
syllables. For $r>17$ all these lists have length at most five,
as required for \cref{cor:twostep}.
\begin{center}\small
\begin{tabular}{r@{\quad}l@{\quad}r}
\toprule
$r$&$(e_1,\ldots,e_L)$&$L$\\
\midrule
7&$(1,-1,1,-1,1,-1,1,-5,-1,1,-1,21,-1,1,-1)$&15\\[2pt]
8&$(1,-1,1,-1,1)$&5\\[2pt]
9&$(1,-1,1,-1,6)$&5\\[2pt]
10&$(1,-1,1,-1,-1,5,-1)$&7\\[2pt]
11&$(1,-1,1,-3,22)$&5\\[2pt]
12&$(1,-1,1)$&3\\[2pt]
13&$(1,-1,1,3,-26)$&5\\[2pt]
14&$(1,-1,1,1,1,2,4,3,-7)$&9\\[2pt]
17&$(15,-1,1,-3,3,-19,4,8,-119)$&9\\[2pt]
24&$(2,-1,2)$&3\\[2pt]
31&$(1,-1,-4,-3,-93)$&5\\[2pt]
35&$(2,-1,5,4,-7)$&5\\[2pt]
37&$(3,-1,3,9,-222)$&5\\[2pt]
59&$(2,-1,-11,9,177)$&5\\[2pt]
61&$(4,-1,7,-89,61)$&5\\[2pt]
83&$(2,-2,13,-15,249)$&5\\[2pt]
85&$(2,-2,15,8,-17)$&5\\[2pt]
107&$(4,-1,-39,49,107)$&5\\[2pt]
109&$(1,-5,10,-50,109)$&5\\[2pt]
613&$(25,-1,-1179,166,1226)$&5\\[2pt]
877&$(2,-18,-135,475,877)$&5\\[2pt]
1139&$(8,-6,701,-617,1139)$&5\\[2pt]
2267&$(1,-94,-206,-1039,-2267)$&5\\[2pt]
\bottomrule
\end{tabular}
\end{center}


\clearpage
\section{Exact matrices for the exceptional words}\label{app:matrices}
The values of the words in \cref{app:words} are
$W_r=\left(\begin{smallmatrix}\alpha&0\\\beta&\alpha^{-1}\end{smallmatrix}\right)$.
The last column is the length of the displayed commutator relation;
no optimality of these lengths is asserted.
\begin{center}
\renewcommand{\arraystretch}{1.4}
\begin{tabular}{r@{\qquad}r@{\qquad}r@{\qquad}r}
\toprule
$r$&$\alpha$&$\beta$&$4L+4$\\
\midrule
7&$\frac{1}{2401}$&$-\frac{410055}{343}$&64\\
8&$1$&$1$&24\\
9&$\frac{1}{9}$&$\frac{2}{3}$&24\\
10&$\frac{1}{25}$&$-\frac{57}{5}$&32\\
11&$\frac{1}{121}$&$\frac{28}{11}$&24\\
12&$-1$&$-1$&16\\
13&$\frac{1}{169}$&$-\frac{46}{13}$&24\\
14&$\frac{1}{2401}$&$-\frac{70501}{343}$&40\\
17&$\frac{1}{83521}$&$-\frac{2444271}{4913}$&40\\
24&$-1$&$-1$&16\\
31&$\frac{1}{961}$&$-\frac{412}{31}$&24\\
35&$\frac{1}{49}$&$-\frac{75}{7}$&24\\
37&$\frac{1}{1369}$&$-\frac{352}{37}$&24\\
59&$\frac{1}{3481}$&$\frac{2848}{59}$&24\\
61&$\frac{1}{3721}$&$\frac{9462}{61}$&24\\
83&$\frac{1}{6889}$&$\frac{7949}{83}$&24\\
85&$\frac{1}{289}$&$-\frac{1050}{17}$&24\\
107&$\frac{1}{11449}$&$\frac{51000}{107}$&24\\
109&$\frac{1}{11881}$&$\frac{54005}{109}$&24\\
613&$\frac{1}{375769}$&$\frac{4798281}{613}$&24\\
877&$\frac{1}{769129}$&$\frac{28102789}{877}$&24\\
1139&$\frac{1}{1297321}$&$\frac{61572851}{1139}$&24\\
2267&$\frac{1}{5139289}$&$-\frac{485429215}{2267}$&24\\
\bottomrule
\end{tabular}
\end{center}


\bibliographystyle{plain}
\bibliography{references}
\end{document}
